\clearpage
\VerdierPhysicalPageAnchor{0069}%
\VerdierSemanticAnchor{ch1-2-2-2-continuation}%
$C$ における $m$-重複複体をなす。

\VerdierSemanticAnchor{ch1-2-2-2-morphism}
次に $f\colon Y\to Y'$ を $C$ における $\varphi$-総和可能な
$n$-重複複体の射とする。次のようにおく：

\VerdierSemanticAnchor{ch1-2-2-2-eq6}
\begin{equation}
  \left(\int_\varphi f\right)^\tau
  =\bigoplus_{\varphi(\sigma)=\tau}f^\sigma\colon
  \left(\int_\varphi Y\right)^\tau
  \longrightarrow
  \left(\int_\varphi Y'\right)^\tau.
  \tag{2.2.2.6}
\end{equation}

\VerdierSemanticAnchor{ch1-2-2-2-functoriality}
$\left(\int_\varphi f\right)^\tau$ が $\int_\varphi Y$ と
$\int_\varphi Y'$ の微分と可換することは明らかである。したがってこの族は
$m$-重複複体の射
$\int_\varphi f\colon\int_\varphi Y\to\int_\varphi Y'$ を定める。また
$\int_\varphi fg=\int_\varphi f\,\int_\varphi g$ および
$\int_\varphi\operatorname{id}_Y=\operatorname{id}_{\int_\varphi Y}$ は
明らかである。ゆえに $\int_\varphi$ は $\varphi$-総和可能な
$n$-重複複体の圏から $m$-重複複体の圏への加法関手である。

\VerdierSemanticAnchor{ch1-2-2-3}
\noindent\textbf{2.2.3.}\quad
$\varphi\colon[n]\to[m]$ を写像、$\tau\in\mathbb Z[m]$、$i\in[m]$ と
する。集合 $E_1$（それぞれ $E_2$）を

\VerdierSemanticAnchor{ch1-2-2-3-fibers}
\[
  \begin{aligned}
    E_1&=\{\sigma_1\in\mathbb Z[n]\mid\varphi(\sigma_1)=\tau\},\\
    \text{（それぞれ }E_2
      &=\{\sigma_2\in\mathbb Z[n]\mid
          \varphi(\sigma_2)=\tau+e_i\}\text{）}
  \end{aligned}
\]
とする。最後に $Y$ を $C$ における $n$-重複複体とし、
$\sigma_1\in E_1$、$\sigma_2\in E_2$ に対して

\VerdierSemanticAnchor{ch1-2-2-3-eq1}
\begin{equation}
  u_{\sigma_2,\sigma_1}(Y)
  =\begin{cases}
    \varepsilon(\tau,i)\varepsilon(\sigma_1,j)d_Y^{\sigma_1,j},
      &\text{もし }\sigma_2-\sigma_1=e_j,\\
    0,&\text{それ以外},
  \end{cases}
  \tag{2.2.3.1}
\end{equation}
とおく。これにより $E_2\times E_1$ 型の行列

\VerdierSemanticAnchor{ch1-2-2-3-eq2}
\begin{equation}
  M(\tau,i,Y)
  =\bigl(u_{\sigma_2,\sigma_1}(Y)\bigr)_{
    \sigma_2\in E_2,\,\sigma_1\in E_1}
  \tag{2.2.3.2}
\end{equation}
を得る。

\VerdierSemanticAnchor{ch1-2-2-3-matrix-identification}
$Y$ が $\varphi$-総和可能ならば、これは射
$d_{\int_\varphi Y}^{\tau,i}$ の行列である。

\VerdierSemanticAnchor{ch1-2-2-3-extensions}
一般に、この行列は各行および各列に有限個しか非零成分をもたない。したがって
$C$ において可算積と直和が表現可能ならば、関手 $\int_\varphi$ を任意の
複体へ拡張する方法は先験的に少なくとも二つある：
